
\subsection*{מטרות השיעור}
\begin{itemize}
    \item חזרה וסיכום על סימטריות ואופרטורים מההרצאה הקודמת
    \item הכרת שלוש הסימטריות הבסיסיות בפיזיקה: $T, C, P$
    \item למידת מושגי יסוד בתורת החבורות (\LR{Groups})
    \item הבנת הצגות (\LR{Representations}) של חבורות ויישומן בפיזיקה
    \item היכרות עם למת שור (\LR{Schur's Lemma}) והשלכותיה
\end{itemize}

\subsection{חזרה וסיכום מהרצאה קודמת}

\subsubsection*{סימטריות ואופרטורים}

\begin{notebox}{תזכורת}
סימטריות במכניקה קוונטית יכולות להיות מיוצגות על ידי:
\begin{itemize}
    \item \textbf{אופרטורים לינאריים ואוניטריים} — רוב הסימטריות (סיבובים, הזזות, שיקופים)
    \item \textbf{אופרטורים אנטי-לינאריים ואנטי-אוניטריים} — היפוך זמן בלבד
\end{itemize}
\end{notebox}

\begin{importantbox}{מסקנה חשובה}
\textbf{כל הסימטריות שאינן כרוכות בהיפוך בזמן הן לינאריות}

אופרטורים אנטי-לינאריים חייבים להיות מצומדים להיפוך בזמן כדי למנוע קריסה למצב מינימום אנרגיה.
\end{importantbox}

\subsection{שלוש סימטריות בסיסיות בפיזיקה}

קיימות שלוש סימטריות בדידות יסודיות בפיזיקה החלקיקים:

\subsubsection*{$T$ (\LR{Time Reversal}) — היפוך בזמן}

\begin{definitionbox}{היפוך זמן}
סימטריה בדידה של היפוך סימן הזמן:
\[
t \to -t
\]

תחת היפוך זמן:
\begin{itemize}
    \item מיקום: $\vec{r}(t) \to \vec{r}(-t)$
    \item מהירות: $\vec{v}(t) \to -\vec{v}(-t)$
    \item תנע: $\vec{p}(t) \to -\vec{p}(-t)$
    \item תנע זוויתי: $\vec{L}(t) \to -\vec{L}(-t)$
\end{itemize}
\end{definitionbox}

\subsubsection*{$C$ (\LR{Charge Conjugation}) — צימוד מטען}

\begin{definitionbox}{צימוד מטען}
סימטריה של היפוך כל המטענים החשמליים:
\[
q \to -q
\]

\textbf{משמעות פיזיקלית:} החלפת חומר באנטי-חומר (אלקטרון $\leftrightarrow$ פוזיטרון)

על פי חוקי ניוטון ומשוואות מקסוול — החוקים הפיזיקליים זהים אם כל המטענים מתהפכים.
\end{definitionbox}

\subsubsection*{$P$ (\LR{Parity}) — שיקוף מרחבי}

\begin{definitionbox}{שיקוף מרחבי (פריטי)}
שיקוף במישור במרחב ("עולם המראה של אליס"):
\[
\vec{r} \to -\vec{r}
\]
\end{definitionbox}

\begin{examplebox}{השפעת שיקוף על גדלים שונים}
    \textbf{\LR{Vectors} (וקטורים רגילים):} מתהפכים בסימן
\begin{itemize}
    \item מיקום: $\pi \circ \vec{r} = -\vec{r}$
    \item מהירות: $\pi \circ \vec{v} = -\vec{v}$
    \item תנע: $\pi \circ \vec{p} = -\vec{p}$
    \item שדה חשמלי: $\pi \circ \vec{E} = -\vec{E}$
\end{itemize}

    \textbf{\LR{Pseudo-vectors} (פסאודו-וקטורים):} נשארים קבועים
\begin{itemize}
    \item תנע זוויתי: $\pi \circ \vec{L} = \vec{L}$ (כי $\vec{L} = \vec{r} \times \vec{p}$)
    \item שדה מגנטי: $\pi \circ \vec{B} = \vec{B}$
\end{itemize}

\textbf{סקלרים:} רוב נשארים קבועים
\begin{itemize}
    \item מסה, טמפרטורה, זמן — קבועים
\end{itemize}

\textbf{פסאודו-סקלרים:} מתהפכים בסימן
\begin{itemize}
    \item נובעים ממכפלה סקלרית של וקטור ופסאודו-וקטור
    \item דוגמה: $\vec{E} \cdot \vec{B}$ (הליציטי)
\end{itemize}
\end{examplebox}

\subsection{שימור סימטריות בכוחות שונים}


\begin{importantbox}{גילוי מרעיש — הפרת הסימטריה}
בשנות ה-50 התגלה כי \textbf{הכוח החלש} (האחראי על דעיכת בטא) \textbf{שובר את הסימטריה} בכל אחד מהצירים בנפרד!

גילוי זה זכה בפרס נובל והפך את ההבנה שלנו על סימטריות בפיזיקה.

\textbf{אבל:} הסימטריה המשולבת \textbf{$CPT$} כן נשמרת!

כלומר: היפוך של זמן + מרחב + מטענים יחד $\Leftarrow$ כוחות מתנהגים זהה.
\end{importantbox}

\subsection{הגדרות בסיסיות בתורת החבורות}

\begin{definitionbox}{חבורה (\LR{Group}) — הגדרה 2.1}
\textbf{חבורה} $G$ היא קבוצה של איברים עם פעולת מכפלה $(\cdot)$ המקיימת ארבעה תנאים:

\begin{enumerate}
    \item \textbf{סגירות (\LR{Closure}):} לכל $a, b \in G$:
    \[
    a \cdot b \in G
    \]
    
    \item \textbf{יחידה (\LR{Identity}):} קיים איבר $e \in G$ כך שלכל $a \in G$:
    \[
    e \cdot a = a = a \cdot e
    \]
    
    \item \textbf{הופכי (\LR{Inverse}):} לכל $a \in G$ קיים $a^{-1} \in G$ כך ש:
    \[
    a \cdot a^{-1} = e = a^{-1} \cdot a
    \]
    
    \item \textbf{אסוציאטיביות (\LR{Associativity}):} לכל $a, b, c \in G$:
    \[
    (a \cdot b) \cdot c = a \cdot (b \cdot c)
    \]

\end{enumerate}
\end{definitionbox}

\begin{notebox}{סימון}
בדרך כלל נשמיט את סימן הפעולה $\cdot$ ונכתוב פשוט $ab$ במקום $a \cdot b$.
\end{notebox}

\subsubsection*{דוגמאות לחבורות}

\begin{examplebox}{דוגמה 1: חבורת החיבור על $\mathbb{R}$}
החבורה $(\mathbb{R}, +)$ — המספרים הממשיים עם פעולת החיבור:
\begin{itemize}
    \item איבר היחידה: $e = 0$
    \item איבר הופכי למספר $a$: $a^{-1} = -a$
    \item חבורה חילופית (אבלית)
\end{itemize}
\end{examplebox}

\begin{examplebox}{דוגמה 2: חבורת סיבובים במישור}
סיבוב בזווית $\theta$ במישור:
\begin{itemize}
    \item איבר היחידה: סיבוב באפס רדיאנים ($\theta = 0$)
    \item איבר הופכי לסיבוב $\theta$: סיבוב $-\theta$
    \item פעולה: $R(\theta_1) \cdot R(\theta_2) = R(\theta_1 + \theta_2)$
    \item חבורה אבלית רציפה
\end{itemize}
\end{examplebox}

\begin{examplebox}{דוגמה 3: חבורת השיקוף}
החבורה $G = \{I, \pi\}$ עם שני איברים:
\begin{itemize}
    \item $I$ — אופרטור הזהות
    \item $\pi$ — אופרטור שיקוף
    \item פעולות: $\pi \circ \pi = I$, $I \circ \pi = \pi$
    \item $\pi$ היא הופכית לעצמה: $\pi^{-1} = \pi$
    \item חבורה בדידה עם 2 איברים
\end{itemize}
\end{examplebox}

\begin{examplebox}{דוגמה 4: החבורה הטריוויאלית}
החבורה $\{I\}$ עם איבר אחד בלבד — איבר היחידה.

זוהי החבורה הפשוטה ביותר (ומשעממת...)
\end{examplebox}

\begin{examplebox}{דוגמה 5: חבורת התמורות $S_3$}
חבורת התמורות של 3 אובייקטים (למשל: 3 כובעים עם 3 כדורים).

\textbf{סימון:} תמורה נכתבת כמטריצה:
\[
\sigma = \begin{pmatrix} 1 & 2 & 3 \\ 2 & 3 & 1 \end{pmatrix}
\]
משמעות: $1 \to 2$, $2 \to 3$, $3 \to 1$

\textbf{גודל החבורה:} $|S_3| = 3! = 6$ איברים

\textbf{תכונה חשובה:} כל חבורה סופית היא תת-חבורה של $S_n$ (משפט קיילי)
\end{examplebox}

\subsection{הצגות של חבורות (\LR{Representations})}

\begin{definitionbox}{הצגה (\LR{Representation}) — הגדרה 2.2}
\textbf{הצגה} של חבורה $G$ היא העתקה (הומומורפיזם):
\[
T: G \to \text{GL}(V)
\]
כאשר $\text{GL}(V)$ היא קבוצת כל האופרטורים הלינאריים ההפיכים על מרחב וקטורי $V$.

\textbf{דרישה:} ההעתקה משמרת את מבנה החבורה:
\[
T(a \cdot b) = T(a) \cdot T(b)
\]
לכל $a, b \in G$.
\end{definitionbox}

\begin{notebox}{פרשנות פיזיקלית}
בפיזיקה קוונטית:
\begin{itemize}
    \item $G$ = חבורת סימטריה (סיבובים, הזזות, וכו')
    \item $V$ = מרחב הילברט של המצבים הקוונטיים
    \item $T(a)$ = אופרטור אוניטרי הפועל על המצבים
\end{itemize}

ההצגה מתארת \textbf{איך הסימטריה משנה את המצבים הקוונטיים}.
\end{notebox}

\subsubsection*{דוגמה: הצגה של $S_3$}

\begin{examplebox}{הצגת חבורת התמורות}
נייצג את $S_3$ על מרחב תלת-ממדי:
\[
\begin{array}{rcl}
1 &\to& (1, 0, 0) \\
2 &\to& (0, 1, 0) \\
3 &\to& (0, 0, 1)
\end{array}
\]

התמורה:
\[
\sigma = \begin{pmatrix} 1 & 2 & 3 \\ 2 & 1 & 3 \end{pmatrix}
\]
מתורגמת למטריצה:
\[
T(\sigma) = \begin{pmatrix} 
0 & 1 & 0 \\ 
1 & 0 & 0 \\ 
0 & 0 & 1 
\end{pmatrix}
\]

כלומר: $(1,0,0) \to (0,1,0)$, $(0,1,0) \to (1,0,0)$, $(0,0,1) \to (0,0,1)$
\end{examplebox}

\subsubsection*{הצגות נאמנות ולא נאמנות}

\begin{examplebox}{הצגות שונות של חבורת השיקוף}
עבור $G = \{I, \pi\}$:

    \textbf{הצגה נאמנה (\LR{faithful}):}
\[
T(I) = 1, \quad T(\pi) = -1
\]
זו הצגה חד-חד ערכית — איברים שונים של החבורה מקבלים אופרטורים שונים.

\textbf{הצגה לא נאמנה:}
\[
T(I) = 1, \quad T(\pi) = 1
\]
ההצגה הטריוויאלית — כל איברי החבורה מקבלים את אופרטור הזהות.
\end{examplebox}

\subsection{משפטים חשובים}

\begin{theorembox}{משפט 2.3 — יוניטריות של הצגות}
לכל חבורה סופית או קומפקטית, כל הצגה שקולה להצגה אוניטרית באמצעות אופרטורים לא סינגולריים.

\textbf{הגדרות:}
\begin{itemize}
    \item \textbf{חבורה סופית:} מספר סופי של איברים
    \item \textbf{חבורה קומפקטית:} "אפשר לשים בקופסא" (בז'רגון פיזיקאי)
    \begin{itemize}
        \item דוגמה: מעגל יחידה $U(1)$ ✓
        \item נגד-דוגמה: הישר הממשי $\mathbb{R}$ ✗
    \end{itemize}
\end{itemize}
\end{theorembox}

\begin{definitionbox}{הצגות שקולות — הגדרה 2.4}
שתי הצגות $T, T'$ של אותה חבורה $G$ נקראות \textbf{שקולות} אם קיימת מטריצה הפיכה $M$ כך ש:
\[
T'(a) = M \, T(a) \, M^{-1}
\]
לכל $a \in G$.

\textbf{משמעות:} שתי הצגות שקולות מתארות את אותה פיזיקה בבסיסים שונים.
\end{definitionbox}

\subsection{דוגמה נגדית: חבורת ההזזות}

\begin{examplebox}{חבורה שאינה קומפקטית}
חבורת ההזזות $T(a): x \to x + a$ בייצוג מטריציוני:
\[
T(a) = \begin{pmatrix} 1 & 0 \\ a & 1 \end{pmatrix}, \quad 
\begin{pmatrix} 1 \\ x \end{pmatrix} \to \begin{pmatrix} 1 \\ x+a \end{pmatrix}
\]

\textbf{בעיה:} המטריצה $T(a)$ \textbf{אינה נורמלית}!
\[
T(a) T(a)^\dagger \neq T(a)^\dagger T(a)
\]

\textbf{הסיבה:} זו חבורה \textbf{לא קומפקטית} ו\textbf{לא סופית} $\Rightarrow$ אינה עומדת בתנאי משפט 2.3
\end{examplebox}

\subsection{פעולה במרחב הפונקציות}

כאשר רוצים להפעיל חבורה על פונקציה $\phi(\vec{x})$, חייבים להיזהר!

\begin{definitionbox}{ההגדרה הנכונה — עם הופכי!}
להפעיל אלמנט $a \in G$ על פונקציה $\phi$:
\[
(a \cdot \phi)(\vec{x}) = \phi(T(a^{-1})\vec{x})
\]

\textbf{שימו לב:} השתמשנו ב-$a^{-1}$ (הופכי) ולא ב-$a$!
\end{definitionbox}

\begin{notebox}{למה דווקא $a^{-1}$?}
כדי לשמור על הפעולה:
\[
(a \cdot b) \cdot \phi = a \cdot (b \cdot \phi)
\]

נבדוק:
\begin{align*}
((a \cdot b) \cdot \phi)(\vec{x}) 
&= \phi(T((ab)^{-1})\vec{x}) \\
&= \phi(T(b^{-1}a^{-1})\vec{x}) \\
&= \phi(T(b^{-1})T(a^{-1})\vec{x})
\end{align*}

מצד שני:
\begin{align*}
(a \cdot (b \cdot \phi))(\vec{x}) 
&= (b \cdot \phi)(T(a^{-1})\vec{x}) \\
&= \phi(T(b^{-1})T(a^{-1})\vec{x})
\end{align*}

זהה! 
\end{notebox}

\begin{warningbox}{ההגדרה השגויה}
אם היינו מגדירים:
\[
(a \cdot \phi)(\vec{x}) = \phi(T(a)\vec{x}) \quad \text{(שגוי!)}
\]

אז:
\[
((a \cdot b) \cdot \phi)(\vec{x}) = \phi(T(ba)\vec{x}) \neq \phi(T(ab)\vec{x})
\]

זה עובד רק לחבורות \textbf{חילופיות} (אבליות)!

לחבורות כלליות — חובה להשתמש בהופכי.
\end{warningbox}

\subsection{מרחב אינווריאנטי}

\begin{definitionbox}{מרחב אינווריאנטי — הגדרה 2.5}
מרחב וקטורי $V$ נקרא \textbf{אינווריאנטי} תחת פעולת חבורה $G$ אם הוא סגור תחת כל האופרטורים:
\[
\forall a \in G, \; \forall \vec{v} \in V: \quad T(a)\vec{v} \in V
\]

\textbf{משמעות:} פעולת החבורה לא "מוציאה" וקטורים מחוץ למרחב.
\end{definitionbox}

\begin{examplebox}{דוגמה}
\textbf{המישור הדו-ממדי $xy$} הוא אינווריאנטי תחת סיבובים \textbf{במישור}, אך \textbf{לא} אינווריאנטי תחת סיבובים כלליים \textbf{במרחב}.
\end{examplebox}

\begin{theorembox}{טענה 2.6 — יצירת מרחב אינווריאנטי}
לכל וקטור $\vec{v}$, המרחב:
\[
V = \text{\LR{Span}}\{T(a)\vec{v} \mid a \in G\}
\]
הוא מרחב אינווריאנטי תחת $G$.

\textbf{הוכחה:} לכל $\vec{w} \in V$ ו-$b \in G$:
\[
\vec{w} = \sum_i c_i T(a_i)\vec{v}
\]
\[
T(b)\vec{w} = \sum_i c_i T(b)T(a_i)\vec{v} = \sum_i c_i T(ba_i)\vec{v} \in V
\]
\end{theorembox}

\subsection{חיבור לפיזיקה — למת שור}

\begin{importantbox}{המשמעות הפיזיקלית}
אם ההמילטוניאן $\hat{H}$ סימטרי תחת חבורה $G$ (כלומר: $[\hat{H}, T(a)] = 0$ לכל $a \in G$), אז:

\textbf{יש ניוון ברמות האנרגיה שקשור להצגות החבורה.}

כל מצב עצמי של האנרגיה שייך להצגה אי-פריקה אחת של החבורה, והניוון מתאים למימד ההצגה.
\end{importantbox}

\begin{notebox}{רקע היסטורי — ישעיהו שור}
    \textbf{ישעיהו שור (\LR{Issai Schur, 1875-1941})} היה מתמטיקאי גרמני-יהודי, תלמיד של פרדינאנד פרוביניוס.

היה פרופסור בברלין עד שגורש בשנת 1935 על פי חוקי הגזע הנאציים. 

ברח לשווייץ ומשם עלה לארץ ישראל, התיישב בתל אביב.

קבור בבית הקברות טרומפלדור בתל אביב.

    \textbf{תרומתו:} למת שור (\LR{Schur's Lemma}) היא אחד הכלים המרכזיים בתורת ההצגות ובפיזיקה קוונטית.
\end{notebox}

\subsection{סיכום עיקרי הנקודות}

\begin{summarybox}{}
\textbf{חזרה על סימטריות:}
\begin{itemize}
    \item סימטריות לינאריות ואוניטריות — רוב הסימטריות הפיזיקליות
    \item סימטריות אנטי-לינאריות — רק היפוך זמן
\end{itemize}

\textbf{שלוש הסימטריות הבסיסיות:}
\begin{itemize}
    \item $T$ — היפוך זמן
    \item $C$ — צימוד מטען (חומר $\leftrightarrow$ אנטי-חומר)
    \item $P$ — שיקוף מרחבי (פריטי)
    \item הכוח החלש שובר כל אחת בנפרד, אך $CPT$ משולב נשמר
\end{itemize}

\textbf{חבורות:}
\begin{itemize}
    \item חבורה = קבוצה + פעולה המקיימת: סגירות, יחידה, הופכי, אסוציאטיביות
    \item דוגמאות: חיבור, סיבובים, תמורות $S_n$
    \item חבורה סופית/קומפקטית: ניתנת להצגה אוניטרית
\end{itemize}

\textbf{הצגות:}
\begin{itemize}
    \item הצגה = הומומורפיזם $T: G \to$ אופרטורים
    \item הצגה נאמנה \LR{vs.} לא נאמנה (טריוויאלית)
    \item פעולה על פונקציות: $(a \cdot \phi)(\vec{x}) = \phi(T(a^{-1})\vec{x})$ — עם הופכי!
\end{itemize}

\textbf{מרחבים אינווריאנטיים:}
\begin{itemize}
    \item מרחב $V$ אינווריאנטי אם $T(a)\vec{v} \in V$ לכל $a \in G$, $\vec{v} \in V$
    \item למה זה חשוב? קשור לניוון ברמות אנרגיה
\end{itemize}

\textbf{קשר לפיזיקה:}
\begin{itemize}
    \item אם $[\hat{H}, T(a)] = 0$ $\Rightarrow$ ניוון שקשור להצגות
    \item למת שור מסבירה את המבנה של מצבים עצמיים
    \item משפט נתר: סימטריה רציפה $\Leftrightarrow$ גודל נשמר
\end{itemize}
\end{summarybox}

\subsection*{נקודות למחשבה ובדיקת הבנה}

\begin{enumerate}
    \item הסבירו מדוע תנע זוויתי $\vec{L} = \vec{r} \times \vec{p}$ הוא פסאודו-וקטור ולא וקטור רגיל. מה קורה לו תחת שיקוף?
    
    \item מדוע הגדרנו את הפעולה על פונקציות עם הופכי: $(a \cdot \phi)(\vec{x}) = \phi(T(a^{-1})\vec{x})$? מה הבעיה אם לא נשתמש בהופכי?
    
    \item תנו דוגמה לחבורה סופית ולחבורה אינסופית קומפקטית. מדוע הישר הממשי $\mathbb{R}$ אינו קומפקטי?
    
    \item הסבירו את ההבדל בין הצגה נאמנה להצגה טריוויאלית. מה המשמעות הפיזיקלית?
    
    \item אם $\hat{H}$ סימטרי תחת חבורה $G$, מה ניתן ללמוד על המצבים העצמיים שלו?
\end{enumerate}

